% C4 — An Amos-type bound at real order, formalized.
% Charter: docs/C4-CHARTER.md (Amendments 1-6).
% REGIME: tricotomy; tag-scoped manuscript hashes; same-lane
% self-citations from v1.  The two former [BIB-ID-SLOT]s were
% WAIVED BY OWNER ORDER at release (charter Amendment 9): the
% companion viXra identifiers are cited as "pending" and will be
% inserted in a post-tag patch when assigned; tags never move.
\documentclass[11pt]{article}
\usepackage[margin=1.1in]{geometry}
\usepackage{amsmath,amssymb,amsthm}
\usepackage{booktabs}
\usepackage[colorlinks=true,linkcolor=blue,urlcolor=blue,citecolor=blue]{hyperref}

\newtheorem{theorem}{Theorem}[section]
\newtheorem{lemma}[theorem]{Lemma}
\newtheorem{remark}[theorem]{Remark}

\newcommand{\lean}[1]{\texttt{#1}}
\newcommand{\RHS}{\mathrm{B}}

\title{A machine-checked proof of an Amos-type bound\\
for modified Bessel ratios at real order}
\author{Lluis Eriksson\thanks{This manuscript and parts of the formal
development were produced with AI assistance under the repository's
audit protocol.  Every theorem below is machine-checked in Lean~4
against a pinned Mathlib; the independent numerical cross-check of
Section~\ref{sec:crosscheck} is a committed interval-arithmetic
transcript; remaining numerics are labelled paper-level where they
occur.  Repository:
\url{https://github.com/lluiseriksson/THE-ERIKSSON-PROGRAMME}.}}
\date{July 2026}

\begin{document}
\maketitle

\begin{abstract}
The Amos-type upper bound for the modified Bessel function ratio,
\[
\rho_\nu(x) \;=\; \frac{I_{\nu+1}(x)}{I_\nu(x)}
\;<\;
\frac{x}{\nu+\tfrac12+\sqrt{(\nu+\tfrac12)^2+x^2}} \;=\; \RHS_\nu(x),
\]
is classical, and its derivation through the qualitative theory of
the associated Riccati equation is an established technique.  A
companion paper formalized the bound at integer order over a
factorial power series.  This paper extends the formalization to
\emph{every real order $\nu \ge 0$}: the function $I_\nu$ is defined
by its $\Gamma$-power series (real exponents via \lean{rpow}), and
the complete chain --- convergence, positivity, the three-term
recurrence, termwise differentiation with a dominated-derivative
argument that must treat the leading term separately (its exponent
$\nu-1$ is negative for $\nu < 1$), the Riccati equation, a
small-argument zone bound uniform in $\nu$, and a first-crossing
barrier --- is machine-checked in Lean~4 with axiom oracle
$[\lean{propext}, \lean{Classical.choice}, \lean{Quot.sound}]$ and
no analytic hypothesis beyond $\nu \ge 0$, $x > 0$.  Two structural
locks tie the result to the integer development: an
\emph{identification theorem} proves that at $\nu = n \in \mathbb{N}$
the $\Gamma$-series object coincides with the factorial-series
object, so the integer-order theorem of the companion development is
recovered as a corollary in three rewrites; and a genuinely
non-integer instance at $\nu = \tfrac12$ witnesses that the endpoint
lives outside the natural-number embedding.  The theorem is proved
for the in-core $\Gamma$-series definition; no identification with
an external special-functions library object is claimed.
\end{abstract}

\section{Introduction}

The ratio $\rho_\nu = I_{\nu+1}/I_\nu$ of modified Bessel functions
of the first kind is bounded above by the algebraically calibrated
expression $\RHS_\nu(x) = x/(\nu+\tfrac12+\sqrt{(\nu+\tfrac12)^2+x^2})$;
bounds of this shape trace back to Amos \cite{Amos74}, and sharp
forms and systematic families were developed by Segura
\cite{Segura11}, Hornik and Gr\"un \cite{HG13,HG24}, and
Ruiz-Antol\'in and Segura \cite{RS16}, the latter describing
explicitly the Riccati-based qualitative methodology that classical
proofs of such bounds employ.  \emph{No novelty is claimed for the
bound or for the Riccati approach.}  We refer to the inequality as
an \emph{Amos-type} bound: the attribution of this exact display to
a specific equation of \cite{Amos74} is deliberately not asserted
here pending collation with the primary source.

The contribution of this paper is the formalization at
\emph{real order}.  Two companion papers of this series formalized,
over the factorial power series at integer order, the consequence
calculus of the bound \cite{ErikssonC2} and then the bound itself
\cite{ErikssonC3}.  The latter states its own scope limit: no
extension to noninteger order is claimed.  This paper removes that
limit.  The bound and its consumers below are stated for every real
$\nu \ge 0$ and every $x > 0$, over the $\Gamma$-power series
\[
\lean{besselIReal}\ \nu\ x \;=\;
\sum_{k \ge 0} \frac{(x/2)^{\nu+2k}}{k!\,\Gamma(\nu+k+1)},
\]
with the real exponent carried by \lean{Real.rpow} and the
denominator by \lean{Real.Gamma}.  In addition, the
\emph{identification theorem} \lean{besselIReal\_natCast} proves
$\lean{besselIReal}\ n = \lean{besselI}\ n$ at every integer order
--- so the integer development of \cite{ErikssonC2,ErikssonC3} is
not a parallel theory but a proved section of the real one, and the
integer-order bound of \cite{ErikssonC3} becomes a corollary
(Theorem~\ref{thm:consumers}).

Statuses: \textbf{exact} means a Lean theorem in the pinned central
repository with clean axiom oracle; \textbf{certified} means a
committed interval-arithmetic transcript; \textbf{paper-level} means
ordinary prose, always flagged; \emph{verified} marks an ordinary
high-precision computation used only as a cross-check of a certified
result, never as evidence.  The repository's largest namespace
(\lean{YangMills/}) reflects its original motivation in lattice
gauge theory; no reduction to any continuum or gauge-theoretic
problem is claimed or used.

\section{Setting}\label{sec:setting}

All objects live in module \lean{AmosClosure} of the pinned central
repository.  From the companion development
\cite{ErikssonC2,ErikssonC3}: the factorial-series
$\lean{besselI}\ n\ x = \sum_{k\ge0}(x/2)^{n+2k}/(k!\,(n+k)!)$ with
its positivity, recurrence, derivative and ratio calculus, the
predicate $\lean{AmosBound}\ \nu\ x\ r \iff r < \RHS_\nu(x)$ ---
already real in its order parameter --- and the calibration
engine.  New in this paper, the real-order object
$\lean{besselIReal}\ \nu\ x$ above, its ratio
$\rho_\nu = \lean{besselRatioReal}\ \nu\ x$, and
$\psi_\nu(x) = x\,(1/\rho_\nu - \rho_\nu)$
(\lean{besselPsiReal}).  Everything is stated at $x > 0$, the
\lean{rpow} domain of the development.

\section{Main results}\label{sec:results}

\begin{theorem}[the Amos-type bound at real order;
\lean{amosBoundReal\_holds}]\label{thm:main}
For every real $\nu \ge 0$ and every real $x > 0$,
\[
\frac{\lean{besselIReal}\ (\nu{+}1)\ x}{\lean{besselIReal}\ \nu\ x}
\;<\; \frac{x}{\nu+\tfrac12+\sqrt{(\nu+\tfrac12)^2+x^2}} .
\]
No hypothesis remains beyond $\nu \ge 0$ and $x > 0$.
\end{theorem}

\begin{theorem}[identification;
\lean{besselIReal\_natCast}]\label{thm:ident}
For every $n \in \mathbb{N}$ and $x > 0$,
\[
\lean{besselIReal}\ (n:\mathbb{R})\ x \;=\; \lean{besselI}\ n\ x .
\]
Termwise, \lean{rpow} collapses to the natural power and
$\Gamma(n+k+1) = (n+k)!$.  The integer-order development is a proved
restriction of the real-order one.
\end{theorem}

\begin{theorem}[consumers, recovery, and the half-order
witness]\label{thm:consumers}
For every real $\nu \ge 0$ and $x > 0$:
$\rho_\nu(x) - \rho_{\nu+1}(x) < 1/x$
(\lean{besselIReal\_unit\_step}); the derivative of
$\log \circ\, \lean{besselIReal}\ \nu$ at $x$ is strictly below that
of $\log \circ\, \lean{besselIReal}\ (\nu{+}1)$, in \lean{deriv}
form (\lean{besselIReal\_logDeriv\_lt}).  Moreover the integer-order
theorem of \cite{ErikssonC3} is recovered through
Theorem~\ref{thm:ident} in three rewrites
(\lean{amosBound\_holds\_recovered}), and the instance $\nu = \tfrac12$
is stated and proved as a named theorem
(\lean{amosBoundReal\_half\_order}) --- the endpoint is not consumed
only through the natural-number embedding.
\end{theorem}

The $\varphi$-monotonicity step of \cite{ErikssonC2} remains
integer-indexed --- its sequence is integer by construction and is
already unconditional through \cite{ErikssonC3}; no artificial
real-order variant is introduced.

\section{The proof}\label{sec:proof}

The formal proofs are the authority; this section presents the
mathematics in full, at the level of the checklist a referee needs.

\subsection{Series facts}\label{sec:series}

Write $t_{\nu,k}(x) = (x/2)^{\nu+2k}/(k!\,\Gamma(\nu+k+1))$.  The
exact term-ratio identity
\[
t_{\nu,k+1} \;=\; t_{\nu,k}\cdot
\frac{(x/2)^2}{(k+1)(\nu+k+1)}
\qquad (\lean{besselRealTerm\_succ})
\]
follows from $\Gamma(s+1) = s\,\Gamma(s)$ at $s = \nu+k+1 \ge 1$ and
$(x/2)^{a+2} = (x/2)^a(x/2)^2$.  Convergence: since
$\Gamma(\nu+1) \le \Gamma(\nu+k+1)$ for $\nu \ge 0$ (each
functional-equation factor $\nu+j \ge 1$), the terms are dominated
by an exponential series, and $I_\nu > 0$ follows termwise.  The
three-term recurrence
$I_\nu - I_{\nu+2} = (2(\nu+1)/x)\,I_{\nu+1}$ is proved termwise
(index shift plus the functional equation), exactly as at integer
order.  Theorem~\ref{thm:ident} is likewise termwise:
$(x/2)^{(n:\mathbb{R})+2k} = (x/2)^{n+2k}$ by
\lean{Real.rpow\_natCast} and
$\Gamma(n+k+1) = (n+k)!$ by \lean{Real.Gamma\_nat\_eq\_factorial}.

\subsection{The derivative layer}\label{sec:deriv}

Termwise differentiation on the interval
$(x/2,\,3x/2) \subset (0,\infty)$ --- the \lean{rpow} domain forbids
a ball through $0$ --- with dominated derivatives.  The termwise
derivative is
\[
\frac{d}{dy}\, t_{\nu,k}(y)
= \frac{\nu+2k}{2}\,
\frac{(y/2)^{\nu+2k-1}}{k!\,\Gamma(\nu+k+1)} .
\]

\begin{remark}[the negative-exponent trap]\label{rem:negexp}
For $k = 0$ and $0 \le \nu < 1$ the exponent $\nu-1$ is negative:
the factor $(y/2)^{\nu-1}$ is \emph{decreasing} in $y$, and a
majorant that evaluates it at the right endpoint is false.  The
formalized domination therefore treats $k=0$ separately, with the
two-endpoint bound
$(y/2)^{\nu-1} \le (x/4)^{\nu-1} + (3x/4)^{\nu-1}$ valid in both
monotonicity regimes (at $\nu = 0$ the whole term carries the factor
$\nu = 0$); and for $k \ge 1$ it uses the factoring
$\nu+2k-1 = (\nu+1) + 2(k-1)$, whose exponents stay $\ge 1$, giving
$(y/2)^{\nu+2k-1} \le (3x/4)^{\nu+1} A^{k-1}$ with $A = (3x/4)^2$.
The majorant tail is summable since
$(\nu+2k+2)/(k+1)! \le (\nu+2)/k!$, equivalent to $0 \le \nu k$.
\end{remark}

Mathlib's dominated-differentiation theorem for series then yields,
for every $\nu \ge 0$ and $x > 0$,
\[
I_\nu'(x) \;=\; I_{\nu+1}(x) + \frac{\nu}{x}\, I_\nu(x)
\qquad (\lean{besselIReal\_hasDerivAt}),
\]
with the \lean{deriv}-form and logarithmic-derivative identities as
corollaries, and the $\nu = n$ value test against the integer layer
through Theorem~\ref{thm:ident}
(\lean{besselIReal\_deriv\_value\_natCast}).

\subsection{The Riccati equation and the nullcline}\label{sec:riccati}

From the derivative identity and the recurrence, the quotient rule
gives
\begin{equation}\label{eq:riccati}
\rho_\nu'(x) \;=\; Q_x(\rho_\nu(x)), \qquad
Q_x(y) := 1 - \frac{2\nu+1}{x}\,y - y^2
\qquad (\lean{besselRatioReal\_hasDerivAt}).
\end{equation}
The bound $\RHS_\nu$ satisfies the exact calibration identity
\begin{equation}\label{eq:calibration}
\frac{1}{\RHS_\nu(x)} - \RHS_\nu(x) \;=\; \frac{2\nu+1}{x}
\qquad (\lean{amosRHS\_calibration\_real}),
\end{equation}
with $a = \nu + \tfrac12 > 0$ supplied by $\nu \ge 0$; multiplied by
$\RHS_\nu > 0$ this is precisely $Q_x(\RHS_\nu(x)) = 0$
(\lean{riccatiQReal\_amosRHS}): the bound is the positive root of
the Riccati quadratic.  Moreover $Q_x(y) > 0$ for
$0 \le y < \RHS_\nu(x)$, by the factorization
$Q_x(y) - Q_x(\RHS_\nu) = (\RHS_\nu - y)\bigl(\tfrac{2\nu+1}{x} + y +
\RHS_\nu\bigr)$ (\lean{riccatiQReal\_pos\_of\_lt}).

\subsection{The zone: $x \le 1/4$, uniformly in $\nu$}\label{sec:zone}

The target $\rho_\nu < \RHS_\nu$ is equivalent, by
\eqref{eq:calibration} and the strict antitonicity of
$t \mapsto 1/t - t$ on $(0,\infty)$, to $\psi_\nu > 2\nu+1$.  On the
zone this is proved directly from the series, with four
product-form facts: geometric domination
$t_{\nu,k} \le t_{\nu,0}\, q_\nu^k$ with $q_\nu = (x/2)^2/(\nu+1)$,
by induction from the term-ratio identity and
$\nu+1 \le (k+1)(\nu+k+1)$
(\lean{besselRealTerm\_le\_geometric}); the product bound
$(1-q_\nu)\, I_\nu \le t_{\nu,0}$ for $q_\nu < 1$
(\lean{besselIReal\_mul\_le});
the strict two-term lower bound $t_{\nu,0} < I_\nu$
(\lean{besselRealTerm\_zero\_lt\_besselIReal}) --- the sole source
of strictness in the whole chain; and the exact identity
$2(\nu+1)\,t_{\nu+1,0} = x\,t_{\nu,0}$, derived by multiplying the
already-proved termwise recurrence at $k=0$ by $x$
(\lean{besselRealTerm\_zero\_succ}).

The $\nu$-uniformity is isolated in two named lemmas rather than
automation: $(\nu+1)/(\nu+2) \le 1$
(\lean{real\_zone\_ratio\_uniform}), whence
$2(\nu+1)\,q_{\nu+1} = \tfrac{x^2}{2}\cdot\tfrac{\nu+1}{\nu+2}
\le \tfrac{x^2}{2}$, and with $x^2 \le 1/16$ the coefficient
inequality
\[
2\nu+1+x^2 \;\le\; 2(\nu+1)\,(1-q_{\nu+1})
\qquad (\lean{real\_zone\_coefficient\_bound}).
\]
Chaining (upper bound at order $\nu+1$, identity, strict lower bound
at order $\nu$) gives the decisive strict inequality
\begin{equation}\label{eq:zonekey}
(2\nu+1+x^2)\; I_{\nu+1}(x) \;<\; x\, I_\nu(x)
\qquad (0 < x \le 1/4),
\end{equation}
strictness entering only at the last link.  Two further lines
convert \eqref{eq:zonekey} into the zone claim: since
$2\nu+1+x^2 \ge 1$, \eqref{eq:zonekey} also gives the crude bound
$I_{\nu+1} \le x\,I_\nu$, hence
$x\,I_{\nu+1}^2 \le x^2\,I_{\nu+1}I_\nu$; multiplying
\eqref{eq:zonekey} by $I_\nu > 0$ and inserting this,
\[
(2\nu+1)\,I_{\nu+1}I_\nu + x\,I_{\nu+1}^2 \;\le\;
(2\nu+1)\,I_{\nu+1}I_\nu + x^2\,I_{\nu+1}I_\nu \;<\; x\,I_\nu^2,
\]
which, by the identity
$\psi_\nu \cdot (I_{\nu+1} I_\nu) = x\,(I_\nu^2 - I_{\nu+1}^2)$, is
precisely $\psi_\nu(x) > 2\nu+1$ on the zone
(\lean{besselPsiReal\_zone}).

\subsection{The barrier: every touch is an upward
crossing}\label{sec:barrier}

Suppose $\psi_\nu(c) = 2\nu+1$ at some $c > 0$.  Then
$1/\rho_\nu(c) - \rho_\nu(c) = (2\nu+1)/c$, and substituting into
the Riccati quadratic annihilates it, so $\rho_\nu'(c) = 0$
(\lean{riccati\_zero\_of\_real\_touch}); the product rule collapses
$\psi_\nu'(c) = (2\nu+1)/c > 0$, the numerator positive directly
from $\nu \ge 0$.  Every touch of the critical level is an upward
crossing.

\begin{lemma}[first crossing; the human form of
\lean{besselPsiReal\_gt}]\label{lem:crossing}
Suppose, for contradiction, $\psi_\nu(x_1) \le 2\nu+1$ for some
$x_1 > 0$.  By the zone theorem $x_1 > 1/4$.  The set
$S = \{y \in [1/4, x_1] : \psi_\nu(y) \le 2\nu+1\}$ is nonempty,
bounded below, and closed (continuity of $\psi_\nu$ on the interval,
from differentiability); let $c = \inf S \in S$.  Since
$\psi_\nu(1/4) > 2\nu+1$, in fact $c > 1/4$, and by minimality
$\psi_\nu > 2\nu+1$ on $(1/4, c)$ (at $1/4$ itself this holds by the
zone theorem); left-continuity forces the touch
$\psi_\nu(c) = 2\nu+1$.  But then $\psi_\nu'(c) = (2\nu+1)/c > 0$,
so the difference quotient
$\bigl(\psi_\nu(y) - \psi_\nu(c)\bigr)/(y-c)$ is eventually positive
as $y \uparrow c$; since $y - c < 0$ there,
$\psi_\nu(y) < \psi_\nu(c) = 2\nu+1$ just left of $c$ ---
contradicting $\psi_\nu > 2\nu+1$ on $(1/4, c)$.  Hence
$\psi_\nu(x) > 2\nu+1$ for all $x > 0$.
\end{lemma}

Finally, $\psi_\nu > 2\nu+1$ converts into $\rho_\nu < \RHS_\nu$
algebraically: with $s = 1/\rho_\nu - \rho_\nu$ and
$t = 1/\RHS_\nu - \RHS_\nu$,
\[
(s - t)\cdot \rho_\nu \RHS_\nu
= (\RHS_\nu - \rho_\nu)\,(1 + \rho_\nu \RHS_\nu),
\]
and the left-hand side is positive --- both its factors are:
$s > t$ is the barrier conclusion $\psi_\nu > 2\nu+1$ divided by
$x$, and $\rho_\nu \RHS_\nu > 0$ --- so, since
$1 + \rho_\nu \RHS_\nu > 0$, positivity passes to
$\RHS_\nu - \rho_\nu$; this is the closing step of
\lean{amosBoundReal\_holds}.  That completes
Theorem~\ref{thm:main}; Theorem~\ref{thm:consumers} is a direct
application, the recovery corollary consisting of the two
identification rewrites, one cast normalization, and \lean{exact}.

\section{Scope}\label{sec:scope}

The theorem is proved for the in-core $\Gamma$-power-series
definition of the modified Bessel function at real order
$\nu \ge 0$, $x > 0$.  No formal identification with an external
special-functions library object is claimed (none exists at the
pinned Mathlib to identify with).  The integer-order objects of the
companion development are identified with the real-order object
\emph{by theorem} (Theorem~\ref{thm:ident}) --- the corresponding
scope caveat of \cite{ErikssonC3} is thereby discharged.  Results
elsewhere in this programme that consume Amos-type bounds for the
classical $I_\nu$ of the analysis literature --- such as the
surface-track $\varphi$-lemma conditional discussed in
\cite{ErikssonC2} --- do not change verification status through
this paper.

\section{Independent numerical cross-check}\label{sec:crosscheck}

A certified interval-arithmetic companion evaluates the margin ---
its protocol pre-registered in the development's charter
(\lean{docs/C4-CHARTER.md}) before the authoritative run ---
$\Delta(\nu,x) = \RHS_\nu(x) - \rho_\nu(x)$ at all $70$ points of
the pre-registered grid
$\nu \in \{0, \tfrac1{10}, \tfrac12, \tfrac9{10}, 1, \tfrac32, 2,
\pi, 10, 100\}$,
$x \in \{\tfrac18, \tfrac14, 1, 2, 5, 50, 300\}$, with $\pi$ a
certified ball.  The ratio is computed by the certified backward
recurrence $\rho_\alpha = x/(2(\alpha+1) + x\rho_{\alpha+1})$,
anchored at order $\nu+N$ with the anchor condition
$q \le \tfrac14$ itself certified in ball arithmetic and seeded by
the two-sided product-form bounds of Section~\ref{sec:zone} --- no
exponentially large value of $I_\nu$ is ever formed, and the seed is
independent of Theorem~\ref{thm:main}.  Verdicts are by ball sign
(\lean{PASS} iff $\inf \Delta > 0$), never by midpoint.  All $70$
points \lean{PASS} at $128$ bits; an independent ordinary
high-precision pass (labelled \emph{verified}, not certified) agrees
at all $70$ points; and at $(\nu,x) = (0,300)$ the certified margin
$[2.78243072419\times10^{-6} \pm 2.43\times10^{-19}]$ reproduces,
by this entirely different method, the slack floor measured by the
integer-order companion of \cite{ErikssonC2}.  It is \emph{not}
part of the proof, and the proof does not cite it.

\section{Related work}\label{sec:related}

Amos \cite{Amos74} introduced bounds of this shape in the context of
computing modified Bessel ratios; two-sided sharp families and
Tur\'an-type consequences are due to Segura \cite{Segura11};
Hornik--Gr\"un studied the Amos-type family systematically
\cite{HG13} and proved in-family optimality results \cite{HG24};
Ruiz-Antol\'in--Segura \cite{RS16} give the sharp bounds at both
parameter ends and describe the Riccati-based qualitative
methodology to which the route formalized here belongs.  The
recurrence and derivative identities are \cite[10.29]{DLMF}, and the
$\Gamma$-series is the standard definition \cite[10.25.2]{DLMF}.  On
the formalization side we know of no prior machine-checked proof of
an Amos-type ratio bound at real order; the integer-order
formalization and its consequence calculus are in the companion
papers \cite{ErikssonC2,ErikssonC3}; the prose companions of the
consumer families are
\cite{ErikssonFH26,ErikssonVM26,ErikssonTuran26}.  The development
builds on Mathlib \cite{mathlib20}.

\section{Reproducibility}\label{sec:repro}

Toolchain: Lean \lean{leanprover/lean4:v4.29.0-rc6}; Mathlib pinned
to commit
\begin{center}\small
\lean{07642720480157414db592fa85b626dafb71355b}
\end{center}
\noindent (lakefile and manifest agree); \lean{tectonic} 0.15.0;
\lean{python-flint} 0.9.0 for the certified companion.

\medskip
\noindent Recorded build facts: at the final mathematics commit
\begin{center}\small
\lean{56ff479b863d8585c41f6738b7bd18f80f2cbb16}
\end{center}
\noindent \lean{lake build AmosClosure} completed successfully
(\textbf{8171 jobs}); the axiom oracle printed, for all \textbf{74}
registered statements (the 34 of the \lean{c3-v1.0} release
registry plus the 40 of this paper), exactly
\begin{center}
\lean{[propext, Classical.choice, Quot.sound]};
\end{center}
\noindent the oracle log is committed under \lean{scripts/} at the
same commit.  The full run history of this development --- eight
build transcripts including the vacuous first run and every failed
attempt, each committed the day it was produced with its diagnosis
in the commit message; the certified-grid transcripts (the
first kept but superseded; the second the authoritative
script-and-transcript pair); and six
charter amendments recording every external audit round and every
method deviation --- is committed under \lean{scripts/} and
\lean{docs/C4-CHARTER.md}; per the repository's audit protocol it is
archival material and plays no role in the mathematics above.

\medskip
\noindent Headline-theorem-to-artifact map (all in module
\lean{AmosClosure}; every row carries a direct oracle witness; the
34 inherited statements are mapped in \cite{ErikssonC2,ErikssonC3}):

\begin{center}
\footnotesize
\begin{tabular}{@{}llll@{}}
\toprule
Result & Lean name & File & Status \\
\midrule
$\Gamma$-series interface & \lean{summable\_besselRealTerm} etc. & \lean{BesselRealInterface.lean} & exact \\
Thm.~\ref{thm:ident} & \lean{besselIReal\_natCast} & \lean{BesselRealInterface.lean} & exact \\
derivative identity & \lean{besselIReal\_hasDerivAt} & \lean{BesselRealDeriv.lean} & exact \\
Eq.~\eqref{eq:riccati} & \lean{besselRatioReal\_hasDerivAt} & \lean{RiccatiReal.lean} & exact \\
Eq.~\eqref{eq:calibration} & \lean{amosRHS\_calibration\_real} & \lean{RiccatiReal.lean} & exact \\
zone (\S\ref{sec:zone}) & \lean{besselPsiReal\_zone} & \lean{AmosBoundProofReal.lean} & exact \\
uniformity lemmas & \lean{real\_zone\_ratio\_uniform} etc. & \lean{AmosBoundProofReal.lean} & exact \\
Lemma~\ref{lem:crossing} & \lean{besselPsiReal\_gt} & \lean{AmosBarrierReal.lean} & exact \\
Thm.~\ref{thm:main} & \lean{amosBoundReal\_holds} & \lean{AmosBarrierReal.lean} & exact \\
consumers & \lean{besselIReal\_unit\_step}, \lean{\_logDeriv\_lt} & \lean{AmosBarrierReal.lean} & exact \\
recovery lock & \lean{amosBound\_holds\_recovered} & \lean{AmosBarrierReal.lean} & exact \\
half-order witness & \lean{amosBoundReal\_half\_order} & \lean{AmosBarrierReal.lean} & exact \\
\S\ref{sec:crosscheck} grid & \lean{c4\_amos\_real\_arb.py} + transcript v2 & \lean{scripts/} & certified \\
\bottomrule
\end{tabular}
\end{center}

\medskip
\noindent From a clean clone at the release revision:
\begin{verbatim}
lake exe cache get
lake build AmosClosure
lake env lean AmosClosure/Oracle.lean
python scripts/c4_amos_real_arb.py
tectonic -X compile papers/c4-amos-real/c4_amos_real.tex
\end{verbatim}
{\sloppy
At the release revision, canonical artifact hashes will be
recorded --- tag-scoped, in \lean{RELEASE-MANIFEST.md} in this
paper's directory --- committed together with the release revision
of this manuscript; compare against
\lean{git show <rev>:path}, never a worktree checkout.  The
annotated tag \lean{c4-v1.0} will mark that release commit after
the pre-release audits close, with the audit registry in
\lean{docs/C4-CHARTER.md}; the release-tree audit is recorded in
the manifest, post-tag by sequencing, with any fixes in post-tag
commits; tags never move.\par}

\section{Conclusion}

An Amos-type bound for modified Bessel ratios is now a theorem of
the pinned formal development at every real order $\nu \ge 0$,
proved from the in-core $\Gamma$-power series using only
machine-checked analytic library results, with no additional
analytic hypothesis beyond $\nu \ge 0$ and $x > 0$, and with
the integer-order development of the companion papers identified as
a proved restriction and recovered as a corollary.  Limitations:
the in-core series definition only, and $\nu \ge 0$
(Section~\ref{sec:scope}); the classical bound and the Riccati
methodology are prior art --- the contribution is the first
machine-checked formalization of this complete chain in the
repository, and we are not aware of a prior machine-checked
formalization of the result.  Natural next steps: the sharp lower
bound (whose small-argument analysis is second-order, as any lower
bound matches the first Taylor coefficient of $\rho_\nu$), formal
identification with a general special-functions interface if one
lands in Mathlib, and upstreaming the modified-Bessel API.

\begin{thebibliography}{99}

\bibitem{Amos74} D.~E.~Amos,
\emph{Computation of modified Bessel functions and their ratios},
Math. Comp. \textbf{28} (1974), 239--251.

\bibitem{Segura11} J.~Segura,
\emph{Bounds for ratios of modified Bessel functions and associated
Tur\'an-type inequalities},
J. Math. Anal. Appl. \textbf{374} (2011), 516--528.
doi:10.1016/j.jmaa.2010.09.030.

\bibitem{HG13} K.~Hornik and B.~Gr\"un,
\emph{Amos-type bounds for modified Bessel function ratios},
J. Math. Anal. Appl. \textbf{408} (2013), 91--101.
doi:10.1016/j.jmaa.2013.05.070.

\bibitem{RS16} D.~Ruiz-Antol\'in and J.~Segura,
\emph{A new type of sharp bounds for ratios of modified Bessel
functions}, J. Math. Anal. Appl. \textbf{443} (2016), 1232--1246.
arXiv:1606.02008.

\bibitem{HG24} K.~Hornik and B.~Gr\"un,
\emph{Generalized Amos-type bounds for modified Bessel function
ratios}, Math. Inequal. Appl. \textbf{27} (2024), 775--787.
doi:10.7153/mia-2024-27-55.

\bibitem{DLMF} NIST Digital Library of Mathematical Functions,
F.~W.~J.~Olver et al., eds., \url{https://dlmf.nist.gov/}.

\bibitem{ErikssonC2} L.~Eriksson,
\emph{One Amos bound, three consumer sites: machine-checked
Bessel-ratio calculus for lattice gauge expansions}, companion
paper, same repository, release \lean{c2-v1.2.1}; ai.viXra.org
identifier pending at the time of this release, July 2026.

\bibitem{ErikssonC3} L.~Eriksson,
\emph{A machine-checked proof of the Amos bound for modified Bessel
function ratios}, companion paper, same repository, release
\lean{c3-v1.0.1}; ai.viXra.org identifier pending at the time of
this release, July 2026.

\bibitem{ErikssonFH26} L.~Eriksson,
\emph{Recurrence--Amos proof of the unit-step order-monotonicity of
$(\log I_\nu)'$, with a Feynman--Hellmann application to
two-dimensional lattice gauge theory}, preprint,
ai.viXra.org:2607.0020, July 2026.

\bibitem{ErikssonVM26} L.~Eriksson,
\emph{Ratio monotonicity for a killed von Mises bridge: exact bridge
structure, certified negative results, a single-Bessel reduction, and
a two-scale closure map for a surface expansion in two-dimensional
lattice gauge theory}, preprint, ai.viXra.org:2607.0023, July 2026.

\bibitem{ErikssonTuran26} L.~Eriksson,
\emph{A weighted Tur\'an-type monotonicity lemma for modified Bessel
functions via the calibrated Amos bound, with an application to the
ordering of a surface expansion in two-dimensional lattice gauge
theory}, preprint, ai.viXra.org:2607.0017, July 2026.

\bibitem{mathlib20} The mathlib Community,
\emph{The Lean mathematical library}, in: Proc. CPP 2020, ACM, 2020,
pp. 367--381. doi:10.1145/3372885.3373824.

\end{thebibliography}

\end{document}
